\documentclass[12pt,a4paper]{article}
\usepackage[utf8]{inputenc}
\usepackage{amsmath}
\usepackage{amsfonts}
\usepackage{amssymb}
\usepackage{graphicx}
\usepackage{natbib}
\usepackage{algorithm}
\usepackage{algpseudocode}

\usepackage{amsthm}
\usepackage{url}
\usepackage{hyperref}
\hypersetup{colorlinks=true,linkcolor=red,citecolor=blue,filecolor=magenta,urlcolor=green}

\title{Seeding Markov Chain Monte Carlo algorithms}
\author{
  Guti\'errez-Pe\~na, Eduardo\\
  \texttt{eduardo@sigma.iimas.unam.mx}
  \and
  Aguirre-P\'erez, Rom\'an\\
  \texttt{roman@sigma.iimas.unam.mx}
}

\begin{document}

\maketitle

\begin{abstract}
The purpose of this work is to illustrate the need of seeding Markov Chain Monte Carlo algorithms when dealing with unobserved variables, as well as to propose a seeding procedure based on a distance-based approach for such instances.  
\end{abstract}

\section{Introduction}
In real-world scenarios, we often aim to explain a set of traits of a phenomenon of interest under specific conditions. From a probabilistic perspective, we usually propose a stochastic mechanism to approximate the underlying law that is governing the phenomenon of interest in such a way that the outcomes of such a stochastic mechanism can be regarded as realisations of the set of traits, say $\pi(\mathbf{z} \vert \boldsymbol{\theta})$, where $\mathbf{z}$ is the set of traits of interest and $\boldsymbol{\theta}$ is a random vector that fully defines the proposed stochastic mechanism.\\

The aforementioned proposal is the core of Statistics, a probabilistic framework for drawing statements about the set of traits by leveraging available measurements of them, say $\mathbf{z}_{1}$, $\mathbf{z}_{2}$, $\ldots$, $\mathbf{z}_{N}$, and denoted by $\mathbf{z}_{1:N}$. In this context, there are two approaches that prevail. On the one hand, the frequentist approach focuses on estimating $\boldsymbol{\theta}$, say $\boldsymbol{\hat{\theta}}$. On the other hand, the Bayesian approach \citep{BS1994} requires the elicitation of $\pi(\boldsymbol{\theta})$, also known as the prior distribution, but returns via Bayes' rule $\pi(\boldsymbol{\theta} \vert \mathbf{z}_{1:N})$, also known as the posterior distribution.\\

We provide a brief overview of the current state-of-the-art MCMC software in Section~\ref{sec:background}. We then address the problem of seeding MCMC algorithms when dealing with unobserved variables and describe the proposed seeding algorithm for overcoming such scenarios in Section~\ref{sec:alg}, whereas in Section~\ref{sec:sim} we illustrate it with a case study of Ecology using the $\mathsf{R}$ packages $\mathbf{NIMBLE}$ \citep{NIMBLE} and $\mathbf{cmdstanr}$ \citep{cmdstanr}. Finally, in Section~\ref{sec:summary} we conclude with....

\section{MCMC in practice}
\label{sec:background}

Markov Chain Monte Carlo (MCMC) methods \citep{HMCMC2011} are widely used in practice for reconstructing posterior distributions. These methods rely on the theoretical principle of defining a homogeneous, Harris recurrent, and aperiodic Markov chain whose stationary probability density function is the target posterior distribution, that is, to construct a suitable Markov chain $\tilde{\boldsymbol{\theta}}^{(1)}, \tilde{\boldsymbol{\theta}}^{(2)}, \ldots$ such that $\tilde{\boldsymbol{\theta}}^{(t)} \underset{t \to \infty}{\sim} \pi(\boldsymbol{\theta} \vert \mathbf{z}_{1:N})$. See \cite{MT2009} for a concise treatment of the theory of Markov chains.\\

Historically, the BUGS project \citep{BUGSbook} led the efforts to enable Bayesian methods  and JAGS \citep{JAGS2003} as one BUGS engine.\\

Lately, Hamiltonian Monte Carlo (HMC) methods \citep{N2011} have became an alternative to MCMC methods.\\

HMC is the workhorse of the $\mathsf{Stan}$ software \citep{Stan2026}. 
On the other hand, the NIMBLE system \citep{NIMBLEpaper}.

\section{Seeding MCMC algorithms}
\label{sec:alg}
 
A more challenging scenario is when we only observed a subset of $\mathbf{z}$, that is, if $\mathbf{z} = (\mathbf{x}, \mathbf{y})$ we only have measurements for $\mathbf{y}$, say $\mathbf{y}^{obs}$. In this case, we instead seek to reconstruct $\pi(\mathbf{x}, \boldsymbol{\theta} \vert \mathbf{y}^{obs})$, and thus, we now also have to initialise $\mathbf{x}$. Built upon a distance-based approach \citep{Maa&etal1996} in tandem with the prior predictive distribution, we proposed the following seeding process for Monte Carlo algorithms when dealing with unobserved variables.\\  


$d:\mathcal{Y}\times\mathcal{Y}\rightarrow\mathbb{R}$ is a pseudo-distance.  Algorithm~\ref{alg:inits} summaries such a seeding process.

\begin{algorithm}
\caption{Distance-based Seeding}
\label{alg:inits}
\begin{algorithmic}[1]

\Require $\boldsymbol{\theta}^{(t)}$
\Ensure $\boldsymbol{\theta}^{(t)}, \mathbf{x}^{(t)}, \mathbf{y}^{(t)}$

\State $t = 0$

\Repeat
	\State $t \leftarrow t+1$
	\State $\boldsymbol{\theta}^{(t)} \sim \pi(\boldsymbol{\theta})$
    \State $\mathbf{x}^{(t)}, \mathbf{y}^{(t)} \sim \pi(\mathbf{x}, \mathbf{y} \vert \boldsymbol{\theta})$
\Until{$d(\mathbf{y}^{(t)}, \mathbf{y}^{obs}) < d^{max}$}

\end{algorithmic}
\end{algorithm}

\section{Case study: Population with exponential growth}
\label{sec:sim}

\citet{AugerMethe2021} provides an overview of Bayesian methods in Ecology, whereas \citet{Monnahan2017} propose to use HMC for making inference.\\

For a discrete-time setting, let us assume that a population experience an exponential growth with time-varying rate $r_{t}$, that is, $N_{t+1}=e^{r_{t}}N_{t}$ for $t=1,2,\ldots,T-1$. We can rewrite the time-varying rate as $r_{t}=f_{t}-m_{t}$, where $f_{t},m_{t}\in\mathbb{R}^{+}$ are the fertility and mortality rates at time $t$, respectively. Now, at each time $t$, we can think of abundance before and after mortality, say $N^{pre}_{t}$ and $N^{post}_{t}$, respectively. Thus, if $M_{t}$ denotes the total mortality we have that $N_{t}^{post}=N_{t}^{pre} - M_{t}$.

Given the aforementioned population demographics, we can propose the following stochastic population mechanism:

\begin{align*}
M_{t} &\sim \mathrm{Bin}\left( N_{t}^{pre}, 1-e^{-m_{t}} \right)\\
N_{t+1}^{pre} &\sim \mathrm{Poi}\left( e^{f_{t}}N_{t}^{post} \right)
\end{align*}

Furthermore,

\begin{align*}
\log(N^{pre}_{1}) &\sim \mathrm{Poi}\left( e^{LN_{0}} \right)\\
\log(a_{t}) &\sim \mathrm{N}\left( \mu_{la}, \sigma_{la} \right)\\
\log(b_{t}) &\sim \mathrm{N}\left( \mu_{lb}, \sigma_{lb} \right)\\
\log(f_{t}) &\sim \mathrm{N}\left( \mu_{lf}, \sigma_{lf} \right)
\end{align*}

with

\begin{align*}
\mu_{lm} &\sim \mathrm{N}\left( m_{lm}, sd_{lm} \right)\\
\sigma_{lm} &\sim \mathrm{N}\left( 0, sd_{lm} \right)\\
\mu_{lf} &\sim \mathrm{N}\left( m_{lf}, sd_{lf} \right)\\
\sigma_{lf} &\sim \mathrm{N}\left( 0, sd_{lf} \right)\\
\end{align*}

The $\mathsf{R}$ project \citep{R2026}.\\

\section{Concluding Remarks}
\label{sec:summary}


\bibliographystyle{apalike}
\bibliography{references.bib}
\cleardoublepage

\end{document}